Observe \(Z=A+\tau G\) with independent standard Gaussian tensor
entries, known \(\tau>0\), norm bound \(\|A\|_F\le M\), and supplied
Tucker or TT caps. Unless specified otherwise, dimensions are equal to
\(p\) and \(K\ge3\). For \(L=p^{-1/2}\sum_{i=1}^p E_i\), write
\(\psi(A)=\langle L,A\rangle\), where \(E_i=e_i^{\otimes K}\). The
criteria \(\mathcal L_0\) and \(\mathcal L_{\max}\) are optimal expected
length at zero and optimal worst-case expected length among intervals
honest over the whole stated class. Fix \(0<\alpha<1/4\).

\subsection{G1. Sharp growing-cap rates with a polynomial
algorithm}\label{g1.-sharp-growing-cap-rates-with-a-polynomial-algorithm}\label{sup:G1}

Let \(r=\min(p,\min_k\bar r_k)\), where the minimum is over the supplied
constrained unfoldings (Tucker modes or TT cuts). Suppose
\(r\ge c\sqrt p\) for a fixed \(c>0\). Uniformly over all remaining
caps, all \(p\), and every \(M,\tau>0\),

\[\mathcal L_0\asymp_{\alpha,c}\mathcal L_{\max}
\asymp_{\alpha,c}\min\{M\sqrt{r/p},\tau\}.\]

Minimax squared risk has the square of this rate. The constants do not
depend on growing ranks or tensor order. Choose deterministically
between the structural interval \([-h_0,h_0]\), \(h_0=M\sqrt{r/p}\), and
the direct Gaussian interval centered at \(\psi(Z)\) with radius
\(z_{1-\alpha/2}\tau\). This takes \(O(p)\) accessed entries and
arithmetic for the coordinate loading. For a rotated flat loading,
obtaining the contractions has its additional polynomial input-access
cost. No factor estimation or exponential scan is required.

\textbf{Upper proof.} An unfolding attaining the smallest cap has
loading matrix singular values \(1/\sqrt p\), repeated \(p\) times.
Matrix singular-value duality and Cauchy--Schwarz bound its contraction
with any rank-at-most-\(r\) matrix of norm \(M\) by \(M\sqrt{r/p}\).
Thus the structural interval covers surely. The direct error is
\(N(0,\tau^2)\). Deterministic radius selection proves the length upper
bound. Choosing estimator zero when \(h_0\le\tau\) and the direct
estimator otherwise gives squared risk at most \(\min(h_0^2,\tau^2)\).

\textbf{Lower proof.} For a uniform size-\(r\) subset \(S\) of
\(\{1,\ldots,p\}\) set

\[A_S=b\sum_{i\in S}E_i,\qquad
b=\min\{M/\sqrt r,\tau\sqrt{\log(1+p/(4r^2))}\}.\]

All constrained unfolding ranks are \(r\), the norm is at most \(M\),
and the common target is \(\Delta=rb/\sqrt p\). For independent supports
\(S,T\), the full Gaussian likelihood second moment equals
\(\mathbb E\exp\{b^2|S\cap T|/\tau^2\}\). To bound it, for
\(H=|S\cap T|\) factorial moments give
\(\mathbb E\binom Hj=\binom rj(r)_j/(p)_j\le\binom rj(r/p)^j\). Expand
\((1+t)^H\) at \(t=e^{b^2/\tau^2}-1\) to obtain

\[1+\chi^2(Q,P_0)\le(1+rt/p)^r\le\exp(r^2t/p)\le e^{1/4}.\]

Put \(x=p/(4r^2)\le1/(4c^2)\). Since \(x/(1+x)\le\log(1+x)\le x\),

\[\frac{\tau}{2\sqrt{1+1/(4c^2)}}
\le\frac{\tau r}{\sqrt p}\sqrt{\log(1+p/(4r^2))}\le\frac\tau2.\]

Consequently \(\Delta\asymp_c\min(M\sqrt{r/p},\tau)\). Total variation
is at most \(\eta=\sqrt{e^{1/4}-1}/2<1/2\). Coverage of both zero and
the common target implies \(\mathbb E_0|C|\ge\Delta(1-2\alpha-\eta)\),
which is positive for \(\alpha<1/4\). The nearest-target testing
argument also gives squared risk at least \(\Delta^2(1-\eta)/8\). This
proves all claims.

This closes the large-cap regime using the earlier sparse-diagonal lower
bound. It is a consequence of that bound with a newly matched upper
bound, not a new sparse-mixture technique. Under \(M=\tau\sqrt p\) the
rate is \(\tau\), so the vanishing fixed-cap improvement has already
disappeared by rank order \(\sqrt p\), well below unrestricted rank
\(p\).

\subsection{G2. General adaptation across positive growing caps is
impossible}\label{g2.-general-adaptation-across-positive-growing-caps-is-impossible}\label{sup:G2}

Set \(M=\tau\sqrt p\) and let the outer smallest cap satisfy
\(r\ge c\sqrt p\). Every interval honest on that outer class has
\(\mathbb E_0|C|\ge c_{\alpha,c}\tau\). For a fixed positive inner cap
vector, the previously established oracle minimax rate is

\[\mathfrak r_K(p)=\frac\tau{\sqrt p}
\begin{cases}(\log(ep))^{3/2},&K=3,\\
\sqrt{\log(ep)},&K\ge4.
\end{cases}\]

Zero belongs to every such inner class. Its worst-case expected length
under outer honesty is therefore at least order \(\tau\), a diverging
factor \(\sqrt p/(\log(ep))^{3/2}\) or \(\sqrt{p/\log(ep)}\) larger than
its known-class oracle rate. No algorithm, including an exponential one,
can be simultaneously minimax in rate over these positive fixed and
growing cap classes. This is stronger than excluding only the
known-rank-zero oracle. It does not conflict with earlier adaptation
across \emph{fixed} outer caps, whose constants were allowed to depend
on that outer cap.

\textbf{Nonzero version.} Suppose \(A_0\) is supported on the first
\(q\) coordinates in every mode, with norm \(a<M\) and actual
constrained ranks \(r_{0,k}\). Let \(m=p-q\),
\(s=\min(m,\min_k(\bar r_k-r_{0,k}))\) and \(M_0=\sqrt{M^2-a^2}\). If
\(s\ge c\sqrt m\), then

\[\mathbb E_{A_0}|C|\gtrsim_{\alpha,c}
\min\{M_0\sqrt{s/p},\tau\sqrt{m/p}\}.\]

\textbf{Proof.} Add the G1 alternatives on the disjoint remaining block.
Norm squares add; every unfolding rank increases by at most \(s\).
Translation by \(A_0\) preserves Gaussian likelihood overlaps, and
normalization of the target multiplies the dimension-\(m\) lower bound
by \(\sqrt{m/p}\). This gives a positive nonzero-base obstruction when
weak extra rank fits inside the supplied caps. Arbitrarily oriented
versions follow by using orthogonal product blocks when the transformed
loading is flat on that block.

\subsection{W1. A dyadic loading bound using the full weight
profile}\label{w1.-a-dyadic-loading-bound-using-the-full-weight-profile}\label{sup:W1}

Return to a prescribed orthogonal product loading
\(L_w=\sum_{i=1}^m w_i U_i\), with decreasing absolute weights and fixed
positive caps. The ambient mode dimensions may exceed \(m\). Partition
the ordered indices into \(B_0=\{1\}\) and
\(B_j=\{2^{j-1}+1,\ldots,\min(2^j,m)\}\) for \(j\ge1\), stopping when
empty. Let \(a_j=\max_{i\in B_j}|w_i|\), \(d_j=|B_j|\), and
\(\alpha_j=\alpha/2^{j+1}\). Define

\[H_K(w)=\sum_j a_j(j+1)^{\nu_K},\qquad
\nu_3=3/2,\quad \nu_K=1/2\ (K\ge4).\]

There is a full-data honest interval satisfying the deterministic bound

\[|C_w|\le C_{K,\bar r,\alpha}
\min\{M\|w\|_\infty,\tau H_K(w)\}.\]

The minimax squared risk is at most a constant times the square of the
same expression. This is a statistical upper bound using the ideal block
residual fits, not a polynomial implementation of those fits.

\textbf{Proof.} Rotate the known orthogonal factors into coordinates.
Restrict the data to each \(d_j^K\) block; this cannot increase its norm
or rank caps. Apply the earlier weighted residual confidence projection
at error \(\alpha_j\). With \(J_2=J(2\bar r)\), its block diameter is at
most

\[4J_2^{K/2}a_j\sum_{s=1}^{d_j}\lambda_s^{(j)}s^{-K/2},\]

where

\[\lambda_s^{(j)}=\tau\sqrt{2\{s\log(ed_j/s)+Ks\log(1+4K)
+\log(2d_j/\alpha_j)\}}.\]

For \(s\ge1\), \(\log(ed_j)\) and \(\log(1/\alpha_j)\) are bounded by a
fixed constant times \(j+1\). Thus
\(\lambda_s^{(j)}\le C_{K,\alpha}\tau\sqrt{s(j+1)}\). The remaining sum
is harmonic for \(K=3\), bounded by \(C(j+1)\), and bounded by a
\(K\)-dependent constant for \(K\ge4\). Add the block intervals and use
\(\sum_j\alpha_j\le\alpha\). Their lengths sum to at most
\(C\tau H_K(w)\). The old structural radius is at most
\(M\sqrt{\min\bar r_k}\|w\|_\infty\), which is \(O(M\|w\|_\infty)\) for
fixed caps. Select the smaller deterministic bound before seeing the
data.

For squared error use block residual fits with error budgets
\(\beta_j=2^{-j-4}\) and deterministic optimization tolerance \(\tau\).
The earlier excess-tail proof gives each block error an \(L_2\) norm at
most \(C\tau a_j(j+1)^{\nu_K}\): the deterministic diameter term has the
bound just derived, and the excess term has \(L_2\) norm \(O(\tau a_j)\)
because \(\sum_s s^{-K/2}<\infty\) and
\(\mathbb E T_{0,j}^2\le2\beta_j\tau^2\). Minkowski's inequality bounds
the \(L_2\) norm of their summed error by \(C\tau H_K(w)\), without
requiring independence. Compare with estimator zero and its structural
risk bound. The finite number of blocks makes measurability and
summation immediate.

\subsection{W2. Sharp rates at every budget for polynomially decreasing
loadings}\label{w2.-sharp-rates-at-every-budget-for-polynomially-decreasing-loadings}\label{sup:W2}

Suppose for fixed \(\beta>0\) and \(C_0<\infty\) that the nonzero,
decreasing weights satisfy

\[|w_i|\le C_0|w_1|i^{-\beta},\qquad i=1,\ldots,m.\]

For fixed positive caps and order, uniformly over all \(m,M,\tau>0\),

\[\mathcal L_0(\Theta;L_w)\asymp\mathcal L_{\max}(\Theta;L_w)
\asymp |w_1|\min(M,\tau),\]

and minimax squared risk is \(\asymp |w_1|^2\min(M,\tau)^2\). Constants
may depend on \(\beta,C_0,K,\bar r,\alpha\), but not on dimension or
budget. This includes all signed power-law weights, whether normalized
or not. In particular, \(0<\beta<1/2\) permits diffuse unit-Frobenius
loadings with \(|w_1|\to0\); the result is not limited to a constant
product overlap with a unit-Frobenius loading.

\textbf{Proof.} For \(j\ge1\), \(a_j\le C_0|w_1|2^{-\beta(j-1)}\). The
series \(\sum_j2^{-\beta(j-1)}(j+1)^{\nu_K}\) is finite for every fixed
\(\beta>0\). W1 therefore gives the matching length and risk upper
bounds. For the lower bounds compare zero with \(tU_1\), where
\(t=\min(M,\tau)/2\) and its sign agrees with \(w_1\). These rank-one
means belong to every positive cap class, their target difference is
\(t|w_1|\), and their total variation is at most \(2\Phi(1/4)-1<1/4\).
The common-target interval argument and two-point squared-loss argument
give the asserted lower orders. No lower power-envelope assumption on
the other weights is needed.

As \(\beta\downarrow0\) the constants deteriorate: the dyadic sum is of
order at most \(1+\beta^{-(\nu_K+1)}\) up to fixed constants. Thus this
theorem does not cover a dimension-dependent exponent approaching zero
uniformly, and does not contradict the extra logarithmic factors for
flat loadings. Completely arbitrary unequal weight sequences, including
many slowly varying profiles, remain uncharacterized.

\textbf{Nonorthogonal perturbations.} Suppose \(L=L_w+R\), with the same
weight condition and a known otherwise arbitrary tensor \(R\) satisfying
\(\|R\|_F\le\epsilon|w_1|\) for a fixed \(0\le\epsilon<1\). The same
sharp rates hold with constants depending also on \(1-\epsilon\). For
the lower bound, \(|\langle L,U_1\rangle|\ge(1-\epsilon)|w_1|\). For
large budgets add a direct interval for \(\langle R,A\rangle\) at a
separate error allocation, costing \(O(\tau\|R\|_F)\). For small
budgets, \(\rho(L_w)=|w_1|\) by generalized Hölder and
\(\rho(R)\le\|R\|_F\); the canonical expansion bounds the structural
radius by \(\sqrt J M(1+\epsilon)|w_1|\). The same reasoning applies to
squared error, using the direct residual estimator and Minkowski. This
gives an all-budget result for a class of loadings without an exact
orthogonal representation, but is not a theorem for every loading.

\subsection{W3. A restricted computational separation for these newly
matched
loadings}\label{w3.-a-restricted-computational-separation-for-these-newly-matched-loadings}\label{sup:W3}

The W2 rate does not automatically have an efficient estimator. To make
this distinction precise, let \(w_i=w_1i^{-\beta}>0\) with fixed
\(0<\beta<1/K\), and set \(M=\tau\sqrt p\). Choose constants

\[\frac1K<a<\min\left\{\frac12,\frac{1-2\beta}{K-2}\right\},
\qquad 0<b<a-\frac1K.\]

For \(s=\lfloor p^a\rfloor\) and \(\mathsf d=\lfloor p^b\rfloor\), every
polynomial estimator of total degree at most \(\mathsf d\) in the full
observed tensor has worst-case squared risk polynomially larger than the
unrestricted \(\tau^2w_1^2\) rate of W2.

\textbf{Finite statement and proof.} For a uniform size-\(s\) support
\(S\), put \(x_S=s^{-1/2}\mathbf1_S\) and let \(\mathsf d_w\) be the
diagonal map with entries \(w_p/w_i\le1\). Use the rank-one alternatives

\[A_S=\tau\lambda(\mathsf d_wx_S)\otimes x_S^{\otimes(K-1)},\qquad
0<\tau\lambda\le M.\]

They have norm at most \(M\) and common weighted target
\(\Delta=\tau\lambda w_p s^{1-K/2}\). All factors are nonnegative, so
for two supports \(0\le\langle A_S,A_T\rangle/\tau^2\le
\lambda^2|S\cap T|^K/s^K\). The full Gaussian Hermite expansion
therefore bounds the squared norm of the centered likelihood ratio
through degree \(\mathsf d\) by

\[\chi^2_{\le \mathsf d}\le\sum_{d=1}^{\mathsf d}
\frac{\lambda^{2d}\mathbb E H^{Kd}}{d!s^{Kd}}
\le\frac{\mu e^\mu B_{\mathsf d}}{1-B_{\mathsf d}}=\eta_{\mathsf d}^2,
\quad \mu=s^2/p,\quad B_{\mathsf d}=\frac{\lambda^2(K\mathsf d)^K}{s^K}<1.\]

For completeness, counting maps by their nonempty fibers gives
\(H^q=\sum_{j=1}^q S(q,j)(H)_j\), where \(S(q,j)\) counts partitions of
\(q\) elements into \(j\) nonempty sets. We have \(S(q,j)\le j^q/j!\)
and \(\mathbb E(H)_j\le\mu^j\). Hence
\(\mathbb E H^q\le q^q(e^\mu-1)\le\mu e^\mu q^q\). Substitute \(q=Kd\)
and sum the geometric upper bound to get the display. For a
degree-\(\mathsf d\) estimator with worst-case squared risk \(R\),
orthogonality of Hermite polynomials and Cauchy--Schwarz give
\(|\mathbb E_QT-\mathbb E_0T|\le\sqrt R\,\eta_{\mathsf d}\). The two
target means differ by \(\Delta\) and their estimation biases are each
at most \(\sqrt R\), so

\[R\ge\Delta^2/(2+\eta_{\mathsf d})^2.\]

Set \(\lambda=\sqrt p\). Then \(\mu\to0\) and
\(B_{\mathsf d}=O(p^{1-Ka+Kb})\to0\), whereas

\[\frac{\Delta}{\tau w_1}
\asymp p^{1/2-\beta-a(K-2)/2}\longrightarrow\infty\]

at a polynomial rate by the choice of \(a\). Such an \(a\) exists
precisely under the stated sufficient condition \(\beta<1/K\). This
proves the restricted separation, using the earlier general low-degree
argument with a new loading-compatible prior. It includes diffuse
power-law loadings with sharp all-budget rates from W2.

This is not a lower bound for all polynomial-time algorithms or all
confidence intervals. An interval honest on zero and these alternatives
would give the test \(\mathbf1\{\Delta\in C(Z)\}\) with total error at
most \(2\alpha+\mathbb E_0|C|/\Delta\). Excluding such efficient tests
requires an additional computational hypothesis; the
polynomial-estimator result alone does not supply one.

\subsection{P1. Efficient pointwise optimality at strong saturated
Tucker
means}\label{p1.-efficient-pointwise-optimality-at-strong-saturated-tucker-means}\label{sup:P1}

This result concerns \textbf{Tucker mode caps all equal to \(r\)},
\(1\le r\le p\). It does not substitute TT bond caps for mode caps. Let

\[A_0=\sum_{i\in S}a_i E_i,\qquad |S|=r,\quad
a=\min_{i\in S}|a_i|>0,\quad
\|A_0\|_F+\tau\le M\le c_0a\sqrt r,\]

where \(c_0>1\) is fixed. The support and signed amplitudes are not
supplied to the algorithm. These means saturate every supplied mode cap.
There is a polynomial interval honest on the entire unseparated Tucker
class such that, in the strong-signal regime specified below,

\[\mathbb E_{A_0}|C|\asymp_{K,\alpha,c_0}\tau\sqrt{r/p},\]

and this matches the infimum expected length at \(A_0\) under that
\emph{same outer-class honesty requirement}. Both \(p\) and \(r\) may
grow. This is pointwise optimality in rate, not optimal leading
constants.

\textbf{Algorithm and noise budgets.} Make the valid artificial split
\(Z^{(1)}=Z+\tau G'\), \(Z^{(2)}=Z-\tau G'\), with independent errors of
scale \(v=\sqrt2\tau\). Set \(\delta=\alpha r/(8p)\) and

\[\eta=2v\sqrt{2\{(p+p^{K-1})\log9+\log(2K/\delta)\}}.\]

Use the preceding cap-valid tangent interval with pilot failure budget
\(\delta\) and total failure budget \(\alpha/2\). Its validation
quantile is \(z_{\mathrm{val}}=z_{1-(\alpha/2-\delta)/2}\). Intersect it
with the direct interval on \(Z\) at error \(\alpha/2\) and the
structural range. All choices are valid on the unseparated class. The
total reported length is bounded by \(2z_{1-\alpha/4}\tau\), including
on pilot failure.

A sufficient condition for the pointwise rate is

\[a\ge 4\eta,\qquad a\ge \eta^2\sqrt p/\tau.\]

For fixed order/confidence this is satisfied at sufficiently large
constant multiple of \(\tau p^{K-1/2}\). This requirement is
conservative and much stronger than a claim of optimal computational
signal strength.

\textbf{Upper proof.} On the simultaneous pilot event, the top-\(r\)
projector \(Q_k\) satisfies \(\|Q_k-P_k\|_{op}\le d=4\eta/a\), where
\(P_k\) projects onto the coordinates in \(S\). The cap-valid
certificate has each angle bound at most \(d\) by Weyl's inequality and
\(a\ge4\eta\). Write \(T_Q\) and \(T_P\) for their first-order
product-space projectors. The representation

\[T_Q=\sum_k\bigotimes_{l\ne k}Q_l-(K-1)\bigotimes_l Q_l\]

with an identity in the omitted mode, and product telescoping, gives
\(\|T_Q-T_P\|_{op}\le2K(K-1)d\). For the diagonal loading,
\(T_PL=p^{-1/2}\sum_{i\in S}E_i\), so

\[\|T_QL\|_F\le\sqrt{r/p}+2K(K-1)d.\]

The cap-valid bias is bounded by \(M\sqrt{\binom K2}d^2\le
16c_0\sqrt{\binom K2}\sqrt r\,\eta^2/a\). The assumed amplitude makes
this \(O(\tau\sqrt{r/p})\). It also gives
\(d\le4\tau/(\eta\sqrt p)=O(\sqrt{r/p})\), because \(\eta\gtrsim\tau\)
and \(r\ge1\). Thus the selected/intersected interval has length
\(O(\tau\sqrt{r/p})\) on the pilot event. Off that event its length is
\(O(\tau)\), with probability at most
\(\delta=O(r/p)\le O(\sqrt{r/p})\). The expected-length upper bound
follows. The dimension-dependent failure budget is essential to this
argument; fixed failure probability would leave an unjustified
\(O(\tau)\) term.

\textbf{Lower proof.} Put \(V=r^{-1/2}\sum_{i\in S}E_i\) and compare
\(A_0\) with \(A_0+(\tau/2)V\). Both have mode ranks at most \(r\) and
norms at most \(M\). Their full Gaussian total variation is less than
\(1/4\) and their target difference is \((\tau/2)\sqrt{r/p}\). The
common-target interval argument at base \(A_0\) proves the lower bound.
This lower comparison uses the same outer class on which the algorithm
is honest.

\textbf{Computation.} Economy SVDs of the \(p\times p^{K-1}\) unfoldings
require \(O(Kp^{K+1})\) arithmetic and \(O(p^K+p^2)\) working storage
for fixed \(K\). Tensor projection has polynomial cost of the same
order. The new code uses economy SVDs; the previous prototype requested
full right singular matrices, which could allocate \(p^{2K-2}\) entries.
Thus the implementation now scales with the explicit tensor rather than
that avoidable allocation. It still reads the full tensor and does not
claim sublinear access cost.

\subsection{What these results do and do not
resolve}\label{what-these-results-do-and-do-not-resolve}

\begin{itemize}
\tightlist
\item
  Efficient high-dimensional minimax computation is resolved here for
  \(r\gtrsim\sqrt p\), where G1 gives the sharp all-budget rate.
  Original fixed-cap diffuse-loading weak signals remain computationally
  open.
\item
  All-budget loading rates are resolved for fixed polynomial decay and
  the specified nonorthogonal perturbations. General profiles remain
  open; the matching W2 statistical construction still uses ideal block
  fits.
\item
  General growing-cap adaptation is impossible in the G2 comparison,
  even between positive cap classes. P1 gives a positive, efficiently
  attained pointwise rate on a strong saturated Tucker family. It does
  not characterize pointwise lengths everywhere or optimal SNR
  thresholds.
\end{itemize}

No blanket tensor hardness theorem is inferred from low-degree results
for sparse tensor PCA. Those results concern a restricted computational
model; a testing reduction for arbitrary efficient intervals still needs
an explicit, justified computational hypothesis.

\subsection{Methodological references}\label{methodological-references}

\begin{itemize}
\tightlist
\item
  The sparse-functional precursor for G1 is Collier, Comminges and
  Tsybakov, \href{https://arxiv.org/abs/1502.00665}{\emph{Minimax
  estimation of linear and quadratic functionals on sparsity classes}}.
\item
  For the distinct estimation-adaptation problem, see Collier,
  Comminges, Tsybakov and Verzélen,
  \href{https://arxiv.org/abs/1611.09744}{\emph{Optimal adaptive
  estimation of linear functionals under sparsity}}.
\item
  The restricted computational perspective in W3 follows the methodology
  of Choo and d'Orsi, \href{https://arxiv.org/abs/2106.06308}{\emph{The
  Complexity of Sparse Tensor PCA}}.
\item
  For general tensor-functional tangent inference under signal
  conditions, see Xu, Chen and Han,
  \href{https://arxiv.org/abs/2501.16223}{\emph{Statistical Inference
  for Low-Rank Tensor Models}}.
\end{itemize}

The repository-specific dependency and novelty distinctions are recorded
in the accompanying audit and literature update.
